\section{Contraction of tree tensor networks}

We first define tree hypergraphs.

\begin{definition}
    % Check with graphical models definitions
    A hypergraph $\graph=(\nodes,\edges)$ is called connected, if for any two nodes $\node,\secnode\in\nodes$ we find a sequence $\nodeof{[n]}$ of nodes, such that $\nodeof{0}=\node,\nodeof{n-1}=\secnode$ and for any $\catenumerator\in[n-1]$ there is $\edgein$ with $\nodeof{\catenumerator},\nodeof{\catenumerator+1}\in\edge$.
    A $\graph=(\nodes,\edges)$ is called a tree, if it is connected and if for any $\edgein$ the hypergraph $(\nodes,\edges\backslash\{\edge\})$ is not connected.
\end{definition}

\subsection{Tree message passing}

\begin{algorithm}[hbt!]
    \caption{Tree Belief Propagation}\label{alg:treeBeliefPropagation}
    \begin{algorithmic}
        \Require Tensor network $\extnet$ on a hypergraph $\graph$ with respect to a semiring $\gensemiring$
        \Ensure Messages $\{\messagewith\wcols(\sedge,\redge)\in\dirovedges\}$
        \iosepline
        \State Initialize a message scheduler $\scheduler=\{(\sedge,\redge) \in \dirovedges \wcols \sedge \text{ a leaf in the tree } (\edges,\dirovedges)\}]$
        \While{$\scheduler$ not empty}
            \State Pop a $(\sedge,\redge)$ pair from $\scheduler$
            \State Compute the message
            \begin{align*}
                \messagewith
                = \srcontractionof{\{\hypercoreofat{\sedge}{\catvariableof{\sedge}}\}
                    \cup \{\mesfromtoat{\secsedge}{\sedge}{\catvariableof{\secsedge\cap\sedge}} \wcols (\secsedge,\sedge)\in\dirovedges \ncond \secsedge\neq \redge\}
                }{\catvariableof{\sedge\cap \redge}}
            \end{align*}
            \State Update $\scheduler$ by all messages $(\redge,\thirdsedge)$ which have not yet been sent, if all messages $(\secsedge,\redge)$ with $\secsedge\neq\thirdsedge$ have been sent.
        \EndWhile
        \State \Return Messages $\{\messagewith\wcols(\sedge,\redge)\in\dirovedges\}$
    \end{algorithmic}
\end{algorithm}


\subsection{Analysis}

We denote for each pair $(\sedge,\redge)$ the subset $\preedgeset\subset\edges$ as the subset of edges $\edgein$, for which each path to $\redge$ passes through $\sedge$.
Note, that by construction $\sedge\in\preedgeset$.

\begin{theorem}
    \label{the:mpGuaranteeTree}
    For any tensor network on a tree hypergraph, \algoref{alg:contractionPropagation} terminates in the tree-based implementation and returns final messages
    \begin{align*}
        \messagewith
        = \srcontractionof{\{\edgehypercorewith\wcols\edge\in\preedgeset\}}{\catvariableof{\sedge\cap\redge}}
    \end{align*}
\end{theorem}
\begin{proof}
    We show this property by induction over the edge sets $\preedgeset$ to pairs $(\sedge,\redge)\in\dirovedges$, such that $\cardof{\preedgeset}\leq n$.
    Notice, that since always $\sedge\in\preedgeset$ we have $n\geq1$.

    $n=1$: In this case we have $\preedgeset=\{\sedge\}$ and $\sedge$ is a leaf of the tree-hypergraph $\graph$.
    The claimed message property holds thus by definition.

    $n\rightarrow n+1$: Let us assume, that the message obeys the claimed property for edge sets with cardinality up to $n$.
    If there is no edge set with cardinality $n+1$, the property holds also for those with cardinality up to $n+1$.
    If there is an edge set $\preedgeset$ with size $n+1$, we have
    \begin{align*}
        \preedgeset = \{\sedge\} \cup \left(\bigcup_{\secsedge \in\dirovedges} \preedgesetwrt{\secsedge}{\sedge}\right) \, .
    \end{align*}
    The message $\mesfromto{\sedge}{\redge}$ is sent, once all messages $\mesfromto{\secsedge}{\sedge}$ to $(\secsedge,\sedge)\in\dirovedges\{(\redge,\sedge)\}$ arrived.
    By definition we have
    \begin{align*}
        \mesfromtowith{\sedge}{\redge}%{\catvariableof{\sedge\cap \redge}}
        = \srcontractionof{\{\hypercoreofat{\sedge}{\catvariableof{\sedge}}\}
            \cup \{\mesfromtowith{\secsedge}{\sedge} \wcols (\secsedge,\sedge) \in\dirovedges \ncond \secsedge \neq \redge \}
        }{\catvariableof{\sedge\cap \redge}}
    \end{align*}
    Now we use the induction assumption on $\preedgesetwrt{\secsedge}{\sedge}$ (since its cardinality is at most $n$) and get
    \begin{align*}
        \mesfromtoat{\sedge}{\redge}{\catvariableof{\sedge\cap \redge}}
        &= \srcontractionof{
            \{\hypercoreofat{\sedge}{\catvariableof{\sedge}}\} \cup
            \left(\bigcup_{(\secsedge,\sedge)\in\dirovedges \ncond \secsedge\neq \redge}
                \srcontractionof{\{\hypercoreofat{\thirdsedge}{\catvariableof{\thirdsedge}} \wcols  \thirdsedge \in \preedgesetwrt{\secsedge}{\sedge}\}}{\catvariableof{\secsedge\cap \sedge}} \right)
        }{\catvariableof{\sedge\cap \redge}} \\
        &= \srcontractionof{
            \{\hypercoreofat{\sedge}{\catvariableof{\sedge}}\} \cup
            \left(\bigcup_{(\secsedge,\sedge)\in\dirovedges \ncond \secsedge\neq \redge} \{\hypercoreofat{\thirdsedge}{\catvariableof{\thirdsedge}} \wcols  \thirdsedge \in \preedgesetwrt{\secsedge}{\sedge}\} \right)
        }{\catvariableof{\sedge\cap \redge}} \\
        &= \srcontractionof{\{\edgehypercorewith\wcols\edge\in\preedgeset\}}{\catvariableof{\sedge\cap\redge}}
    \end{align*}
    Here we used the commutation of contraction property in the second equation, which is justified by the assumed tree property of the hypergraph.
    Thus, the message property holds also for any edge sets of size $n+1$.

    By induction, the claimed message property therefore holds for all final messages.
\end{proof}
